% This file is part of the modular source for Stratified Relational Workspace Theory.
% !TeX root = ../Stratified_Relational_Workspace_Theory.tex
\theorypart{II}{Static Core--Transfer Theory}

\topic{3}{Interaction Structure}

\begin{definition}[Interaction graph]\label{def:interaction-graph}
Let
\[
G=(V,E,s,t)
\]
be a finite directed multigraph. Parallel edges and self-loops are allowed.
The vertex set indexes strata, participants, artifacts, models, instruments,
or subworkspaces. An edge \(e:i\to j\) denotes an admissible directed
translation or influence channel from \(i\) to \(j\).
\end{definition}

The directed path category is
\[
\mathcal C=\Path(G).
\]
Its objects are vertices. A path
\[
\gamma=e_m\cdots e_2e_1
\]
means that \(e_1\) is traversed first. Composition is concatenation, and the
empty path at \(i\) is denoted \(1_i\).

For every \(i\in V\), assign a finite-dimensional Hilbert space
\[
X_i.
\]
For every edge, assign a linear transfer
\[
T_e:X_{s(e)}\longrightarrow X_{t(e)}.
\]
The path transfer is
\[
T_\gamma=T_{e_m}\cdots T_{e_2}T_{e_1},
\qquad
T_{1_i}=I_{X_i}.
\]
Accordingly, the static transfer system is a functor
\[
\mathcal X:\Path(G)\longrightarrow\mathbf{FinHilb}.
\]

The finite-dimensional Hilbert assumption supplies adjoints and spectral
tools. The exact compatible-section definitions require less and may later be
formulated in categories such as
\[
\mathbf{Ban},\qquad
\mathbf{Meas},\qquad
\mathbf{Prob},\qquad
\mathbf{Top},\qquad
\mathbf{Alg}.
\]

\subtopic{3.1 Directed paths versus invertible transport}

Global compatible sections make sense even when transfers are noninvertible.
Holonomy, however, requires transport around reversible loops. To state the
exact holonomy theorem without ambiguity, let \(|G|\) be the underlying
undirected graph and let
\[
\Pi_1(G)
\]
denote its fundamental path groupoid. Every oriented edge is then supplied
with a formal reverse. The core transfers extend to \(\Pi_1(G)\) only when
the relevant maps \(H_e\) are invertible, with
\[
H_{e^{-1}}=H_e^{-1}.
\]

Hence two regimes will be distinguished:
\begin{enumerate}
\item the \emph{diagram regime}, in which arbitrary directed maps define
compatible sections; and
\item the \emph{local-system regime}, in which invertible core maps define a
functor
\[
H:\Pi_1(G)\to\mathbf{FinHilb}^{\times}.
\]
\end{enumerate}
Theorems based on holonomy will explicitly assume the second regime.

\subtopic{3.2 Connectivity convention}

Unless stated otherwise, connectedness means that \(|G|\) is connected. A
root \(o\in V\) may then be joined to every vertex by a path in the underlying
graph. If \(G\) has several connected components, the theory applies to each
component and
\[
\Gamma(H)
\cong
\bigoplus_{\alpha}\Gamma(H|_{G_\alpha}).
\]

\subtopic{3.3 Edge weights}

For approximate theory, assign strictly positive edge weights
\[
w_e>0.
\]
They may encode channel reliability, frequency, precision, experimental
importance, or inverse noise variance. Exact compatibility is independent of
positive rescaling of these weights, while approximate spectra are not.

\begin{example}[Human--model--artifact graph]\label{ex:hma-graph}
The motivating three-node system is represented by
\[
V=\{h,m,a\},
\qquad
E=\{h\to m,m\to a,a\to h\}.
\]
Additional edges may encode direct retrieval \(a\to m\), human revision
\(h\to a\), or independent verification. The graph records permitted
translations, not merely physical communication links.
\end{example}

\topic{4}{Core--Residual Decomposition}

\begin{definition}[Admissible local decomposition]\label{def:core-residual}
For every vertex \(i\), choose an orthogonal decomposition
\[
X_i=C_i\oplus R_i,
\]
with projectors
\[
P_i:X_i\to C_i,
\qquad
Q_i=I-P_i:X_i\to R_i.
\]
The subspace \(C_i\) is the candidate locally realized transferable core. The
subspace \(R_i\) contains substrate-specific, context-specific, or otherwise
noncore degrees of freedom relative to the present analysis.
\end{definition}

The decomposition is part of the model and is not presumed unique. Calling a
subspace ``core'' is therefore a hypothesis requiring structural and causal
validation. Later null tests must determine whether the chosen projectors
identify a genuine shared channel rather than an arbitrary low-dimensional
summary.

Relative to these decompositions, every edge map has the block form
\[
T_e=
\begin{pmatrix}
H_e & U_e\\
V_e & K_e
\end{pmatrix},
\]
where
\[
H_e=P_{t(e)}T_eP_{s(e)}:C_{s(e)}\to C_{t(e)},
\]
\[
U_e=P_{t(e)}T_eQ_{s(e)}:R_{s(e)}\to C_{t(e)},
\]
\[
V_e=Q_{t(e)}T_eP_{s(e)}:C_{s(e)}\to R_{t(e)},
\]
\[
K_e=Q_{t(e)}T_eQ_{s(e)}:R_{s(e)}\to R_{t(e)}.
\]

\subtopic{4.1 Structural conditions}

\textbf{Core preservation} means
\[
V_e=0.
\]
A core input cannot leak into the target residual.

\textbf{Core autonomy} means
\[
U_e=V_e=0.
\]
Core and residual evolve as separate transfer systems.

\textbf{Core leakage}, \textbf{residual contamination}, and \textbf{total
transfer defect} are measured by
\[
\lambda_e^{\mathrm{out}}=\|V_e\|,
\qquad
\lambda_e^{\mathrm{in}}=\|U_e\|,
\]
\[
\delta_e^{\mathrm{tr}}
=
\|U_e\|+\|V_e\|.
\]
These quantities depend on the chosen norms and projectors. Their empirical
meaning must therefore be tied to a declared intervention and prediction
protocol.

\subtopic{4.2 Composition exposes residual mediation}

For composable edges \(e:i\to j\) and \(f:j\to k\), direct block
multiplication gives
\[
T_fT_e
=
\begin{pmatrix}
H_fH_e+U_fV_e & H_fU_e+U_fK_e\\
V_fH_e+K_fV_e & V_fU_e+K_fK_e
\end{pmatrix}.
\]
The induced core-to-core block is therefore
\[
P_kT_fT_eP_i
=
H_fH_e+U_fV_e.
\]
The term \(U_fV_e\) is a residual-mediated correction: information leaves the
candidate core at the intermediate node and later re-enters it. Consequently,
the simple path product \(H_fH_e\) faithfully represents full core transport
only under a preservation condition such as \(V_e=0\), or under an explicit
model that includes these mediated corrections.

\begin{proposition}[Pathwise core preservation]\label{prop:pathwise-preservation}
Suppose \(V_e=0\) for every edge in a directed path
\(\gamma=e_m\cdots e_1\). Then
\[
T_\gamma(C_{s(\gamma)})\subseteq C_{t(\gamma)}
\]
and
\[
P_{t(\gamma)}T_\gamma P_{s(\gamma)}
=
H_{e_m}\cdots H_{e_1}.
\]
\end{proposition}

\begin{proof}
For a single edge, \(V_e=0\) means
\[
T_e(c,0)=(H_ec,0),
\]
so core states remain in the core. Inductively, if the first \(r\) transfers
map \(c\) to \((H_{e_r}\cdots H_{e_1}c,0)\), applying \(T_{e_{r+1}}\) gives
\[
(H_{e_{r+1}}H_{e_r}\cdots H_{e_1}c,0).
\]
The conclusion follows after \(m\) steps.
\end{proof}

\begin{corollary}[Autonomous splitting]\label{cor:autonomous-splitting}
If \(U_e=V_e=0\) for every edge, then the complete diagram decomposes as the
direct sum of two diagrams,
\[
\mathcal X\cong H\oplus K,
\]
and every path transfer satisfies
\[
T_\gamma
=
\begin{pmatrix}
H_\gamma&0\\
0&K_\gamma
\end{pmatrix}.
\]
\end{corollary}

\subtopic{4.3 Exact core-transfer systems}

\begin{definition}[Exact core-transfer system]\label{def:exact-core-transfer}
An exact core-transfer system on \(G\) consists of:
\begin{enumerate}
\item decompositions \(X_i=C_i\oplus R_i\);
\item core-preserving edge maps \(V_e=0\); and
\item invertible core blocks \(H_e\in GL(C_{s(e)},C_{t(e)})\) whenever
holonomy is invoked.
\end{enumerate}
\end{definition}

The requirement \(V_e=0\) concerns the relation between the selected core and
the full transfer. The requirement that \(H_e\) be invertible concerns the
ability to compare core representations along reversed paths. They answer
different questions and should not be conflated.

\begin{remark}[Approximate systems]
In applications neither condition need hold exactly. The approximate theory
will replace strict invariance by low transfer energy, small leakage, and a
spectral gap separating workspace modes from residual modes.
\end{remark}

\topic{5}{Abstract Operation Algebra}

\begin{definition}[Operation algebra]\label{def:operation-algebra}
Let \(\mathbb F\in\{\mathbb R,\mathbb C\}\). An operation algebra
\(\mathcal A\) is a unital associative \(\mathbb F\)-algebra whose elements
represent abstract transformations, inferential rules, questions, or
procedures that can be realized in multiple strata.
\end{definition}

Every core space carries a unital representation
\[
\rho_i:\mathcal A\longrightarrow\End(C_i).
\]
For \(a\in\mathcal A\), the operator \(\rho_i(a)\) is the realization of the
same abstract operation in local coordinates.

\begin{definition}[Exact operation-compatible transfer]\label{def:intertwiner}
A core transfer \(H_e:C_i\to C_j\) is operation compatible when it is an
intertwiner:
\[
\boxed{
H_e\rho_i(a)=\rho_j(a)H_e
}
\qquad
\forall a\in\mathcal A.
\]
Equivalently, the diagram
\[
\begin{array}{ccc}
C_i&\xrightarrow{\rho_i(a)}&C_i\\
\downarrow H_e&&\downarrow H_e\\
C_j&\xrightarrow{\rho_j(a)}&C_j
\end{array}
\]
commutes for every \(a\).
\end{definition}

This condition states that translating and then operating gives the same
result as operating and then translating. Agreement of static encodings alone
is weaker: two nodes may store corresponding vectors while implementing
incompatible transformations on them.

\begin{proposition}[Path composition of intertwiners]\label{prop:path-intertwiner}
If every edge of a path \(\gamma:i\to j\) intertwines the local
\(\mathcal A\)-representations, then
\[
H_\gamma\rho_i(a)=\rho_j(a)H_\gamma
\qquad
\forall a\in\mathcal A.
\]
\end{proposition}

\begin{proof}
For composable intertwiners \(H_e:C_i\to C_j\) and
\(H_f:C_j\to C_k\),
\[
H_fH_e\rho_i(a)
=H_f\rho_j(a)H_e
=\rho_k(a)H_fH_e.
\]
Induction over path length proves the claim.
\end{proof}

\begin{corollary}[Holonomy commutes with abstract operations]\label{cor:holonomy-commutes}
For every loop \(\gamma:o\to o\),
\[
H_\gamma\rho_o(a)=\rho_o(a)H_\gamma.
\]
Hence
\[
\rho_o(\mathcal A)
\subseteq
\operatorname{Comm}(\Hol_o),
\]
where \(\operatorname{Comm}(\Hol_o)\) is the commutant of the holonomy
representation.
\end{corollary}

\subtopic{5.1 Approximate operation compatibility}

The edgewise defect is
\[
\Xi_e(a)
=
H_e\rho_{s(e)}(a)-\rho_{t(e)}(a)H_e.
\]
For a bounded generating family \(\mathcal G\subseteq\mathcal A\), define
\[
\eta_{\mathrm{int}}(\mathcal G)
=
\sup_{a\in\mathcal G}
\sum_{e\in E}w_e\|\Xi_e(a)\|^2.
\]
When \(\mathcal A\) is normed and its unit ball is appropriate, one may take
\[
\mathcal G=\{a:\|a\|_{\mathcal A}\leq1\}.
\]
Exact representational universality means
\[
\eta_{\mathrm{int}}=0,
\]
while approximate universality means
\[
\eta_{\mathrm{int}}\leq\varepsilon_{\mathrm{int}}.
\]
The use of a generating family prevents the definition from depending on an
unbounded rescaling of algebra elements.

\begin{example}[Equivalent operations in different coordinates]\label{ex:coordinate-intertwiner}
Let \(C_1=C_2=\mathbb R^2\), choose \(S\in GL_2(\mathbb R)\), and set
\[
H_{12}=S,
\qquad
\rho_2(a)=S\rho_1(a)S^{-1}.
\]
Then
\[
H_{12}\rho_1(a)
=S\rho_1(a)
=\rho_2(a)S
=\rho_2(a)H_{12}.
\]
The local matrices can look different while realizing the same abstract
operation. This is the elementary model of representation change used
throughout the theory.
\end{example}
